\documentclass[11pt,a4paper]{article}
\usepackage[utf8]{inputenc}
\usepackage[T1]{fontenc}
\usepackage[french]{babel}
\usepackage{geometry}
\usepackage{enumitem}
\usepackage{parskip}
\usepackage{microtype}

\geometry{margin=2cm}
\setlist[itemize]{noitemsep, topsep=2pt, parsep=0pt, partopsep=0pt}
\setlength{\parskip}{0.4em}

\title{Rapport d'avancement}
\author{Application Pacing (desktop)}

\begin{document}
\maketitle
\thispagestyle{empty}
\vspace{-1.5em}

%----------------------------------------------------------------------
\section{Reformulation des missions et tâches confiées}
%----------------------------------------------------------------------

Les travaux confiés visent à \textbf{intégrer les données USA Swimming} dans l'application existante (déjà alimentée par les données françaises), puis à \textbf{exploiter ces données dans le couloir de performance} de la même manière que pour la France. Ils comprennent également \textbf{l'intégration des données des nageurs marocains}, afin de pouvoir \textbf{comparer leurs performances} avec celles des nageurs français et américains sur les mêmes graphiques de couloir.

%----------------------------------------------------------------------
\section{Actions réalisées depuis le 15 mai}
%----------------------------------------------------------------------

\begin{itemize}
  \item Branchement des données USA dans l'application desktop (onglet dédié, aperçu des jeux de données), pour vérifier que le chargement se fait sans problème.
  \item Mise en cache et chargement structuré des compétitions USA Swimming.
  \item Ajout du couloir global pour les données américaines, en complément de celui déjà en place pour la France, avec une liste déroulante «\,Pays\,» pour basculer entre les deux jeux de données.
  \item Mise en place d'une barre de recherche des nageurs dans le couloir de performance d'USA Swimming, pour choisir le nageur de référence.
  \item Sélection du nageur marocain sur les couloirs France\,/\,États-Unis (liste déroulante et barre de recherche), pour superposer ses performances sur le graphique.
\end{itemize}

\noindent\textit{Travaux connexes (hors liste principale)\,:} nettoyage et préparation des données USA, notebooks d'exploration et de contrôle qualité, harmonisation des données marocaines (FRM Natation) pour la comparaison.

%----------------------------------------------------------------------
\section{Difficultés et points de blocage rencontrés}
%----------------------------------------------------------------------

\begin{itemize}
  \item \textbf{Volume et lenteur des données USA} : Premier chargement et passage France $\rightarrow$ États-Unis très longs.
  \item \textbf{Recalcul du couloir à chaque changement d'épreuve} : Interface qui se figeait lors des changements de distance, nage ou bassin.
  \item \textbf{Recherche et sélection de nageurs} : Suggestions trop lourdes, bouton de confirmation parfois inopérant.
  \item \textbf{Qualité des noms (USA)} : Noms invalides (chiffres, caractères spéciaux) perturbant listes et recherche.
  \item \textbf{Extraction des résultats marocains (PDF)} : Les PDF issus du scraping du site FRM Natation présentaient des mises en page et structures très hétérogènes selon les compétitions\,; la structuration du contenu a nécessité un choix de modèles d'IA adapté à la taille du document (nombre de tokens), pour rester dans les limites des modèles tout en obtenant des données exploitables.
  \item \textbf{Comparaison Maroc\,/\,France\,/\,USA} : Alignement des épreuves, tranches d'âge et formats de temps entre sources différentes.
  \item \textbf{Complexité de l'interface} : gestion simultanée de plusieurs contrôles (nageur marocain, bouton de confirmation et indicateur de chargement).


\end{itemize}

%----------------------------------------------------------------------
\section{Solutions mises en œuvre pour les dépasser}
%----------------------------------------------------------------------

\textbf{Données USA}
\begin{itemize}
  \item Stockage des compétitions en \textbf{Parquet par année} et \textbf{chargement parallèle} de plusieurs années.
  \item \textbf{Cache en mémoire} après le premier chargement pour limiter les relectures complètes.
  \item \textbf{Filtrage des noms invalides} lors de la préparation des données.
\end{itemize}

\textbf{Couloir de performance et interface}
\begin{itemize}
  \item \textbf{Cache en mémoire des images} de couloir déjà calculées\,; réaffichage prioritaire lors d'un changement d'épreuve si une image existe déjà, recalcul en arrière-plan sinon.
  \item \textbf{Calcul du graphique en arrière-plan} pour ne plus bloquer l'interface.
  \item \textbf{Recherche avec attente courte après la saisie} (debounce), suggestions et confirmation du nageur\,; indicateur de chargement pendant la recherche.
  \item \textbf{Limitation des listes de nageurs} affichées pour éviter la création de milliers d'éléments à l'écran.
\end{itemize}

\textbf{Données marocaines (extraction et comparaison)}
\begin{itemize}
  \item Utilisation du service \textbf{LlamaIndex} pour l'\textbf{extraction} et la \textbf{structuration} du contenu des fichiers PDF issus du scraping FRM Natation, avec production de données conformes à un \textbf{schéma JSON défini} (compétitions, épreuves, performances, nageurs).
  \item Pipeline de nettoyage FRM Natation, tranches d'âge et nages alignées sur le modèle France\,/\,USA.
  \item Chargeur dédié avec \textbf{cache en mémoire} après le premier chargement.
\end{itemize}

%----------------------------------------------------------------------
\section{Éléments qui fonctionnent correctement et avancées obtenues}
%----------------------------------------------------------------------

\textbf{Fonctionnel aujourd'hui}
\begin{itemize}
  \item Chargement et aperçu des données USA via l'onglet de test dédié.
  \item Cache Parquet par année et chargement parallèle opérationnels.
  \item Couloir global \textbf{France} et \textbf{États-Unis} accessibles via la liste «\,Pays\,».
  \item Recherche et sélection d'un nageur de référence sur le couloir (France\,/\,USA).
  \item Superposition d'un nageur marocain sur les couloirs France et États-Unis (liste + recherche).
  \item Temps de réponse du couloir \textbf{nettement améliorés} après optimisation.
\end{itemize}

\textbf{Avancées obtenues}
\begin{itemize}
  \item Les données USA sont \textbf{intégrées et exploitables} dans la même logique que les données françaises pour le couloir global.
  \item L'utilisateur peut \textbf{comparer visuellement} un profil marocain à un couloir de référence français ou américain sur une même épreuve.
  \item La base technique (Parquet, parallélisme, cache mémoire) est en place pour \textbf{absorber le volume} des compétitions américaines.
\end{itemize}


\vfill

\end{document}
